\section*{Bab \ref{ch:g9:arith} --- Aritmetika: Pembagi dan Bilangan Prima}

\begin{solution}{exo:g9:arith:1}
\omterm{def:g9:arith:divisor}{Pembagi} $36$: $1, 2, 3, 4, 6, 9, 12, 18, 36$.
\omterm{def:g9:arith:divisor}{Pembagi} $45$: $1, 3, 5, 9, 15, 45$.
\omterm{def:g9:arith:divisor}{Pembagi} $17$: hanya $1$ dan $17$ ($17$ adalah prima).
\end{solution}

\begin{solution}{exo:g9:arith:2}
$2\,346$ berakhir dengan $6$: \omterm{def:g6:wholes:divisible}{habis dibagi} $2$, tidak \omterm{def:g6:wholes:divisible}{habis dibagi}
$5$. \omterm{def:g1:addition:def}{Jumlah} angkanya $2 + 3 + 4 + 6 = 15$: \omterm{def:g6:wholes:divisible}{habis dibagi} $3$, tidak
\omterm{def:g6:wholes:divisible}{habis dibagi} $9$. Dua angka terakhirnya
$46$, dan $46 = 4 \times 11 + 2$ tidak \omterm{def:g6:wholes:divisible}{habis dibagi} $4$: jadi
$2\,346$ tidak \omterm{def:g6:wholes:divisible}{habis dibagi} $4$.
\end{solution}

\begin{solution}{exo:g9:arith:3}
$72 = 2^3 \times 3^2$; \quad
$150 = 2 \times 3 \times 5^2$; \quad
$210 = 2 \times 3 \times 5 \times 7$; \quad
$121 = 11^2$.
\end{solution}

\begin{solution}{exo:g9:arith:4}
$101$: ujilah \omterm{def:g9:arith:prime}{bilangan prima} $p$ dengan $p^2 \leq 101$, yaitu
$2, 3, 5, 7$. Tidak ada yang membagi $101$ (\omterm{def:g3:numbers:evenodd}{ganjil}, \omterm{def:g1:addition:def}{jumlah} angkanya
$2$, tidak berakhir dengan $0/5$, dan
$101 = 7 \times 14 + 3$): jadi $101$ prima.

$91 = 7 \times 13$: bukan prima.

$143 = 11 \times 13$: bukan prima.
\end{solution}

\begin{solution}{exo:g9:arith:5}
\omterm{def:g9:arith:divisor}{Pembagi} bersama $48$ dan $60$: \omterm{def:g9:arith:divisor}{pembagi} 48 adalah $1, 2, 3, 4, 6, 8,
12, 16, 24, 48$; \omterm{def:g9:arith:divisor}{pembagi} $60$ adalah $1, 2, 3, 4, 5, 6, 10, 12, 15,
20, 30, 60$; yang bersamanya adalah $1, 2, 3, 4, 6, 12$, jadi
$\gcd(48,60) = 12$.

Dengan pemfaktoran: $48 = 2^4 \times 3$ dan
$60 = 2^2 \times 3 \times 5$; \omterm{def:g9:arith:prime}{bilangan prima} bersamanya dengan
\omterm{def:g8:powers:def}{eksponen} yang lebih kecil: $2^2 \times 3 = 12$.
\end{solution}

\begin{solution}{exo:g9:arith:6}
$\gcd(255, 154)$:
\begin{align*}
255 &= 154 \times 1 + 101, \\
154 &= 101 \times 1 + 53, \\
101 &= 53 \times 1 + 48, \\
53 &= 48 \times 1 + 5, \\
48 &= 5 \times 9 + 3, \\
5 &= 3 \times 1 + 2, \\
3 &= 2 \times 1 + 1, \\
2 &= 1 \times 2 + 0 .
\end{align*}
\omterm{def:g3:division:remainder}{Sisa} terakhir yang bukan \omterm{def:g1:counting:zero}{nol}: $\gcd(255, 154) = 1$ (keduanya \omterm{def:g9:arith:gcd}{saling prima}).

$\gcd(1053, 325)$:
\begin{align*}
1053 &= 325 \times 3 + 78, \\
325 &= 78 \times 4 + 13, \\
78 &= 13 \times 6 + 0 .
\end{align*}
$\gcd(1053, 325) = 13$.
\end{solution}

\begin{solution}{exo:g9:arith:7}
Algoritma Euclid: $588 = 504 \times 1 + 84$; $504 = 84 \times 6 + 0$:
jadi $\gcd(588, 504) = 84$. Lalu
\[
\frac{588}{504} = \frac{588 \div 84}{504 \div 84} = \frac{7}{6},
\]
yang tak tersederhanakan.
\end{solution}

\begin{solution}{exo:g9:arith:8}
Banyaknya rangkaian harus membagi $84$ dan $126$; yang terbesar
mungkin adalah $\gcd(84, 126)$. Pemfaktorannya: $84 = 2^2 \times 3
\times 7$, $126 = 2 \times 3^2 \times 7$, jadi FPB-nya
$2 \times 3 \times 7 = 42$. Ia dapat membuat $42$ rangkaian, yang
masing-masing berisi
$\frac{84}{42} = 2$ mawar dan $\frac{126}{42} = 3$ tulip.
\end{solution}

\begin{solution}{exo:g9:arith:9}
Kita perlu \omterm{def:g9:arith:divisor}{kelipatan} bersama terkecilnya. $24 = 2^3 \times 3$ dan
$36 = 2^2 \times 3^2$; dengan mengambil tiap \omterm{def:g9:arith:prime}{bilangan prima} beserta
\omterm{def:g8:powers:def}{eksponen} yang \emph{lebih besar}: $\lcm = 2^3 \times 3^2 = 72$.
Ferinya berikutnya berangkat bersama $72$ menit sesudah pukul 8:00,
yaitu pukul 9:12.
\end{solution}

\begin{solution}{exo:g9:arith:10}
\emph{1.} Setiap \omterm{def:g9:arith:divisor}{pembagi} bersama $d$ dari $n$ dan $n+1$ juga membagi
\omterm{ex:g1:subtraction:difference}{selisihnya} $(n+1) - n = 1$, jadi $d = 1$: yaitu $\gcd(n, n+1) = 1$.

\emph{2.} Sebuah \omterm{def:g9:fractions:fraction}{pecahan} tak tersederhanakan persis ketika \omterm{def:g4:fractions:def}{pembilang}
dan \omterm{def:g4:fractions:def}{penyebutnya} \omterm{def:g9:arith:gcd}{saling prima}, dan itulah keadaan $n$ dan $n + 1$
menurut bagian 1.
\end{solution}

\begin{solution}{pb:g9:arith:1}
\textbf{1.} Isilah kendi $3$ lalu tuangkan ke kendi $5$ (isinya
$0/3 \to 3$ pada yang besar). Isilah kendi $3$ lagi lalu tuangkan ke
kendi $5$ sampai penuh: kendi besarnya hanya menerima $2$ lagi,
sehingga meninggalkan
\[
3 - 2 = 1 \text{ L pada kendi kecilnya.}
\]
Gerakannya: isi $3$; tuang $3 \to 5$; isi $3$; tuang $3 \to 5$.

\textbf{2.} Isilah kendi $5$; tuangkan ke kendi $3$ (sehingga
tersisa $2$ pada yang besar); kosongkan kendi $3$; tuangkan $2$-nya
ke kendi $3$; isilah kendi $5$; tuangkan ke kendi $3$ sampai penuh
--- kendi itu menerima $1$, sehingga meninggalkan
$\mathbf{4}$ L pada kendi besarnya. Enam gerakan.

\textbf{3.} Semuanya: $1$ (pertanyaan 1), $2$ (sesudah dua gerakan
pada pertanyaan 2), $3$ dan $5$ (sekali pengisian), $4$ (pertanyaan
2), $6 = 3 + 3$ (kendi kecil yang penuh ditambah $3$ yang dituang ke
yang besar), $7 = 5 + 2$, $8 = 5 + 3$ (keduanya penuh). Setiap takaran
bulat dari $1$ sampai $8$ L dapat ditakar dengan kendi $5$ dan $3$.

\textbf{4.} Awalnya: kedua kendinya memuat $0$, yaitu \omterm{def:g9:arith:divisor}{kelipatan} $2$.
Mengisi menetapkan isinya menjadi $6$ atau $4$: \omterm{def:g3:numbers:evenodd}{genap}. Mengosongkan
menetapkannya menjadi $0$: \omterm{def:g3:numbers:evenodd}{genap}. Menuang memindahkan sebagian airnya
antara kendi yang isinya \omterm{def:g3:numbers:evenodd}{genap}, dan banyaknya yang dituang adalah
\omterm{ex:g1:subtraction:difference}{selisih} \omterm{def:g3:numbers:evenodd}{bilangan genap} (ruang yang tersisa, atau banyaknya yang
tersedia): jadi semua isinya tetap \omterm{def:g3:numbers:evenodd}{genap} selamanya. Sasaran \omterm{def:g3:numbers:evenodd}{ganjil}
seperti $1$ L tidak tercapai.

\textbf{5.} $\gcd(6, 4) = 2$: jadi hanya takaran \omterm{def:g3:numbers:evenodd}{genap} --- yang
dibenarkan pertanyaan 4. $\gcd(5, 3) = 1$: setiap takaran bulat
diizinkan hukumnya, dan pertanyaan 3 mewujudkan semuanya. FPB itu
persis satuan ukur kedua kendinya.

\textbf{6.} $91 = 1 \times 65 + 26$; $65 = 2 \times 26 + 13$;
$26 = 2 \times 13 + 0$: jadi $\gcd(91, 65) = 13$. Dan
$2\,026 = 44 \times 46 + 2$; $46 = 23 \times 2 + 0$:
jadi $\gcd(2\,026, 46) = 2$.

\textbf{7.} Menuang kendi $5$ ke kendi $13$ berulang-ulang: sesudah
dua kali pengisian kendi besarnya memuat $10$; pengisian ketiganya
hanya muat $3$, sehingga meninggalkan $5 - 3 = 2$ pada kendi kecilnya
--- \omterm{def:g3:division:remainder}{sisa}
$13$ dibagi $5$ adalah $3$, dan takaran $3$ (ruangnya) dan $2$
(kelebihannya) persis bilangan milik Euclid ($13 = 2 \times 5 + 3$,
$5 = 1 \times 3 + 2$). Kalau dilanjutkan, $3 - 2 = 1$ muncul: yaitu
\omterm{def:g3:division:remainder}{sisa} berikutnya pada algoritmanya. Pancurannya menjalankan pembagian
Euclid dengan air.

\textbf{8.} $\gcd(13, 5) = 1$, jadi hukum pada pertanyaan 5
mengizinkan setiap takaran bulat --- dan air terjun pada pertanyaan 7
memang menghasilkan $1$ L. Ya, dapat.

\textbf{9.} \omterm{def:g9:arith:divisor}{Pembagi} bersama $n$ dan $2n + 1$ membagi
$2n + 1 - 2 \times n = 1$: jadi ia haruslah $1$. Karena itu
$\gcd(n, 2n+1) = 1$ selalu.

\textbf{10.} Bus A: 7:12, 7:24, 7:36, 7:48, 8:00\,\dots; bus B:
7:18, 7:36, 7:54\,\dots\ Keberangkatan bersama yang pertama: 7:36,
sesudah $36$ menit --- yaitu \omterm{def:g9:arith:divisor}{kelipatan} bersama pertama $12$ dan $18$.
Hukumnya: $36 \times \gcd(12, 18) = 36 \times 6 = 216 = 12 \times
18$. Untuk $5$ dan $3$: \omterm{def:g9:arith:divisor}{kelipatan} bersama pertamanya $15$, dan
$15 \times \gcd(5,3) = 15 \times 1 = 15 = 5 \times 3$.

\textbf{11.} Dengan daur $17$ tahun: keberbarengan berikutnya adalah
\omterm{def:g9:arith:divisor}{kelipatan} bersama pertama $17$ dan $4$; karena
$\gcd(17, 4) = 1$, itu berarti $17 \times 4 = 68$ tahun --- jangkriknya
bertemu puncaknya sekali dalam empat kemunculan. Dengan daur
$16$ tahun: $16$ adalah \omterm{def:g9:arith:divisor}{kelipatan} $4$, jadi \emph{setiap}
kemunculannya mengenai sebuah puncak. Panjang daur yang prima tidak
berbagi faktor dengan daur pemangsa yang lebih pendek mana pun,
sehingga mendorong keberbarengannya sejauh mungkin: aritmetika sebagai
penyamaran.

\textbf{12.} Empat pilihan \omterm{def:g8:powers:def}{eksponen} untuk $2$, tiga untuk $3$,
dua untuk $5$: jadi $4 \times 3 \times 2 = 24$ \omterm{def:g9:arith:divisor}{pembagi}.

\textbf{13.} Kalau $m = 2^{a} \times 3^{b} \times \cdots$, maka
$m^2 = 2^{2a} \times 3^{2b} \times \cdots$: setiap \omterm{def:g8:powers:def}{eksponennya}
berlipat dua, jadi \omterm{def:g3:numbers:evenodd}{genap}. Pada $360 = 2^3 \times 3^2 \times 5$,
\omterm{def:g8:powers:def}{eksponen} $2$ dan \omterm{def:g8:powers:def}{eksponen} $5$ \omterm{def:g3:numbers:evenodd}{ganjil}: jadi $360$ bukan kuadrat
sempurna.

\textbf{14.} Sebuah \omterm{def:g9:arith:divisor}{pembagi} menjadi pasangannya sendiri persis ketika
$d = \frac nd$, yaitu $n = d^2$: hanya kuadrat yang punya \omterm{def:g9:arith:divisor}{pembagi}
tengah seperti itu. Untuk semua $n$ yang lain \omterm{def:g9:arith:divisor}{pembaginya} terbelah
menjadi pasangan, jadi banyaknya \omterm{def:g3:numbers:evenodd}{genap}. Jadi: banyak \omterm{def:g9:arith:divisor}{pembagi} yang
\omterm{def:g3:numbers:evenodd}{ganjil} $\Leftrightarrow$ kuadrat sempurna. Periksa: $36$ punya \omterm{def:g9:arith:divisor}{pembagi}
$1, 2, 3, 4, 6, 9, 12, 18, 36$ --- sembilan buah, yaitu \omterm{def:g3:numbers:evenodd}{ganjil}, dan
$36 = 6^2$; sedangkan $360$ punya $24$ (pertanyaan 12), yaitu \omterm{def:g3:numbers:evenodd}{genap},
dan bukan kuadrat (pertanyaan 13).

\textbf{15.} Loker $n$ terbalik keadaannya sekali oleh tiap murid $k$
yang nomornya membagi $n$: seluruhnya sebanyak \omterm{def:g9:arith:divisor}{pembagi} yang dipunyai
$n$. Sebuah loker berakhir \emph{terbuka} ketika keadaannya terbalik
sebanyak \omterm{def:g3:numbers:evenodd}{bilangan ganjil} kali --- menurut pertanyaan 14, persis ketika
$n$ adalah kuadrat sempurna. Loker yang terbuka:
$1, 4, 9, 16, 25, 36, 49, 64, 81, 100$ --- sepuluh buah.

\end{solution}
